\section*{الفصل \ref{ch:g2:caring-for-nature} --- العناية بالطبيعة}

\begin{solution}{exo:g2:caring-for-nature:1}
ما يحيط \omterm{def:g1:living-or-not:living}{بالكائن الحي} ويعتمد عليه. مثلًا: الهواء الذي نتنفسه،
والماء الذي نشربه، والغذاء الذي نأكله --- وكلها من \omterm{def:g2:caring-for-nature:environment}{الطبيعة}.
\end{solution}

\begin{solution}{exo:g2:caring-for-nature:2}
مثلًا: الخنافس تحت اللحاء، وحشرات القمل الخشبي تحت الجذع، والطحلب
فوقه، والقنفذ الذي يشتّي في جوفه.
\end{solution}

\begin{solution}{exo:g2:caring-for-nature:3}
مثلًا: احمل كل النفايات إلى البيت؛ وانظر بدل أن تجمع، ولا تقطف
حزمة كاملة أبدًا؛ وأعد كل حجر أو جذع تقلبه؛ والزم الممرات قرب
المواضع الهشّة؛ وراقب \omterm{def:g1:animals-around-us:animal}{الحيوانات} في هدوء من بعيد ودع صغارها وشأنها.
\end{solution}

\begin{solution}{exo:g2:caring-for-nature:4}
لأن \omterm{def:g2:caring-for-nature:environment}{الطبيعة} لا شاحنة نفايات لها: فما نتركه يبقى --- والغطاء
البلاستيكي يبقى سنوات كثيرة كثيرة، تكفي لأن يوقع \omterm{def:g1:animals-around-us:animal}{الحيوانات} في شرك
أو يسمّمها، وليبقى ملقًى هناك حين يكبر أطفال اليوم.
\end{solution}

\begin{solution}{exo:g2:caring-for-nature:5}
إن قُطفت الحزمة إلى آخرها لم تستطع أن تصنع بذورًا --- فلا أقحوان
هناك في السنة القادمة --- ويفقد النحل الأزهار التي يتغذى منها.
\end{solution}

\begin{solution}{exo:g2:caring-for-nature:6}
تتعفن ببطء فتصير تربة داكنة هشّة. والديدان وحشرات القمل الخشبي
وكائنات حية أصغر من أن تُرى تقوم بالعمل؛ وننال في النهاية سمادًا
--- تربة جيدة تغذّي \omterm{def:g1:plants-around-us:plant}{نباتات} السنة القادمة.
\end{solution}

\begin{solution}{exo:g2:caring-for-nature:7}
مثلًا: صحن ماء في موجة حرّ، وركن بري في الحديقة يُترك بلا جزّ،
\omterm{def:g2:from-seed-to-plant:seed}{وبذور} للطيور في البرد، وحفر بركة صغيرة.
\end{solution}

\begin{solution}{exo:g2:caring-for-nature:8}
الجواب مفتوح. والأغلفة البلاستيكية والقناني تفوز عادة. وكل صنف كان
ينبغي أن يعود مع صاحبه إلى البيت --- ثم إلى حاوية الفرز الصحيحة.
\end{solution}

\begin{solution}{exo:g2:caring-for-nature:9}
اتركي \omterm{ex:g2:animals-grow-up:frog}{الشراغيف} في البركة: ففيها كل ما تحتاج إليه --- ماء البركة،
والغذاء، ومتسع \omterm{def:g1:plants-around-us:plant}{لنبات} \omterm{def:g1:my-body:parts}{الأرجل} --- وليس ذلك في المرطبان. وإن أرادت
إيما أن تساعد فلتراقبها هناك طول الربيع، ولتُبقِ الضفة غير مدوسة،
ولعلها تعين على صنع \omterm{def:g2:where-animals-live:habitat}{موئل} بركة جديد كبركة المدرسة.
\end{solution}

\begin{solution}{exo:g2:caring-for-nature:10}
مليون ``لا شيء'' في الاتجاه السيئ مرج مدفون تحت الأغطية ---
فالأضرار الصغيرة تتضاعف. وفي الاتجاه الحسن، مليون عناية صغيرة ---
نفايات محمولة إلى البيت، وصنابير مغلقة، ونفايات مفروزة --- تتراكم
بالقدر نفسه من اليقين: \omterm{def:g2:caring-for-nature:environment}{فالطبيعة} يُعتنى بها في خيارات بهذا الصغر
تمامًا.
\end{solution}
